\documentclass[12pt]{article}
\usepackage[a4paper,margin=1in]{geometry}
\usepackage[utf8]{inputenc}
\usepackage[T1]{fontenc}
\usepackage{amsmath,amssymb}
\usepackage{textcomp}
\usepackage{parskip}
\begin{document}
\begin{center}
{\Large \textbf{2010 STEP III}}
\end{center}

\subsection*{Page 1}
\begin{flushleft}
\bigskip\textbf{Section A: Pure Mathematics}\par
\bigskip\textbf{Question 1.} Let x , x , . . . , x and x be any fixed real numbers. The numbers A and B are defined by\par
1 2 n n+1\par
n n\par
1 1\par
\[A = x , B = (x - A) ,\]
n k n k\par
\[Xk=1 Xk=1\]
and the numbers C and D are defined by\par
n+1 n+1\par
1 1\par
\[C = x , D = (x - C) .\]
n + 1 k n + 1 k\par
\[Xk=1 Xk=1\]
\textbf{(i)} Express C in terms of A, x and n.\par
n+1\par
n\par
2 2\par
\textbf{(ii)} Show that B = x - A .\par
n k\par
\[Xk=1\]
\textbf{(iii)} Express D in terms of B, A, x and n.\par
n+1\par
\[Hence show that (n + 1)D > nB for all values of x , but that D < B if and only if\]
n+1\par
(n + 1)B (n + 1)B\par
\[A - < x < A + .\]
r n n+1 r n\par
\bigskip\textbf{Question 2.} In this question, a is a positive constant.\par
\textbf{(i)} Express cosh a in terms of exponentials.\par
By using partial fractions, prove that\par
1 a\par
\[dx = .\]
\[Z^0 x^2 + 2x cosh a + 1 2 sinh a\]
\textbf{(ii)} Find, expressing your answers in terms of hyperbolic functions,\par
\infty\par
dx\par
\[Z^1 x^2 + 2x sinh a - 1\]
and\par
\infty\par
dx .\par
\[Z^0 x^4 + 2x2 cosh a + 1\]
\[9475 Jun10\]
\end{flushleft}
\subsection*{Page 2}
\begin{flushleft}
\bigskip\textbf{Question 3.} For any given positive integer n, a number a (which may be complex) is said to be a primitive\par
nth root of unity if an = 1 and there is no integer m such that 0 < m < n and am = 1. Write\par
down the two primitive 4th roots of unity.\par
Let C (x) be the polynomial such that the roots of the equation C (x) = 0 are the primitive\par
n n\par
nth roots of unity, the coefficient of the highest power of x is one and the equation has no\par
\[repeated roots. Show that C^4(x) = x^2 + 1 .\]
\textbf{(i)} Find C1(x), C2(x), C3(x), C5(x) and C6(x), giving your answers as unfactorised poly-\par
nomials.\par
\textbf{(ii)} Find the value of n for which C (x) = x4 + 1.\par
n\par
\textbf{(iii)} Given that p is prime, find an expression for C (x), giving your answer as an unfactorised\par
p\par
polynomial.\par
\textbf{(iv)} Prove that there are no positive integers q, r and s such that C (x) \equiv C (x)C (x) .\par
q r s\par
\bigskip\textbf{Question 4.} (i) The number α is a common root of the equations x2 + ax + b = 0 and x2 + cx + d = 0\par
\[(that is, α satisfies both equations). Given that a 6= c, show that\]
b - d\par
\[α = - .\]
a - c\par
Hence, or otherwise, show that the equations have at least one common root if and only\par
if\par
2 2\par
\[(b - d) - a(b - d)(a - c) + b(a - c) = 0 .\]
\[Does this result still hold if the condition a 6= c is not imposed?\]
\textbf{(ii)} Show that the equations x2 + ax + b = 0 and x3 + (a + 1)x2 + qx + r = 0 have at least\par
one common root if and only if\par
2 2\par
\[(b - r) - a(b - r)(a + b - q) + b(a + b - q) = 0 .\]
\[Hence, or otherwise, find the values of b for which the equations 2x2 + 5x + 2b = 0 and\]
\[2x3 + 7x2 + 5x + 1 = 0 have at least one common root.\]
\[9475 Jun10 [Turn over\]
\end{flushleft}
\subsection*{Page 3}
\begin{flushleft}
\bigskip\textbf{Question 5.} The vertices A, B, C and D of a square have coordinates (0, 0), (a, 0), (a, a) and (0, a),\par
respectively. The points P and Q have coordinates (an, 0) and (0, am) respectively, where\par
\[0 < m < n < 1. The line CP produced meets DA produced at R and the line CQ produced\]
meets BA produced at S. The line P Q produced meets the line RS produced at T . Show\par
that T A is perpendicular to AC.\par
Explain how, given a square of area a2, a square of area 2a2 may be constructed using only a\par
straight-edge.\par
[Note: a straight-edge is a ruler with no markings on it; no measurements (and no use of\par
compasses) are allowed in the construction.]\par
\bigskip\textbf{Question 6.} The points P , Q and R lie on a sphere of unit radius centred at the origin, O, which is fixed.\par
\[Initially, P is at P^0(1, 0, 0), Q is at Q^0(0, 1, 0) and R is at R^0(0, 0, 1).\]
\textbf{(i)} The sphere is then rotated about the z-axis, so that the line OP turns directly towards\par
the positive y-axis through an angle φ. The position of P after this rotation is denoted\par
\[by P^1. Write down the coordinates of P^1.\]
\textbf{(ii)} The sphere is now rotated about the line in the x-y plane perpendicular to OP1, so that\par
the line OP turns directly towards the positive z-axis through an angle λ. The position\par
\[of P after this rotation is denoted by P^2. Find the coordinates of P^2. Find also the\]
\[coordinates of the points Q^2 and R^2, which are the positions of Q and R after the two\]
rotations.\par
\textbf{(iii)} The sphere is now rotated for a third time, so that P returns from P2 to its original\par
\[position P^0. During the rotation, P remains in the plane containing P^0, P^2 and O. Show\]
that the angle of this rotation, θ, satisfies\par
\[cos θ = cos φ cos λ ,\]
and find a vector in the direction of the axis about which this rotation takes place.\par
\[9475 Jun10\]
\end{flushleft}
\subsection*{Page 4}
\begin{flushleft}
\bigskip\textbf{Question 7.} Given that y = cos(m arcsin x), for |x| < 1, prove that\par
\[d2y dy\]
2 2\par
\[(1 - x ) - x + m y = 0 .\]
\[dx2 dx\]
\[d3y d2y dy\]
Obtain a similar equation relating , and , and a similar equation relating\par
\[dx3 dx2 dx\]
\[d4y d3y d2y\]
, and .\par
\[dx4 dx3 dx2\]
dn+2y dn+1y dny\par
Conjecture and prove a relation between , and .\par
dxn+2 dxn+1 dxn\par
Obtain the first three non-zero terms of the Maclaurin series for y. Show that, if m is an even\par
integer, cos mθ may be written as a polynomial in sin θ beginning\par
\[m^2 sin 2 θ m^2(m^2 - 22) sin 4 θ\]
\[1 - + - \cdots . (|θ| < \pi)\]
2! 4! 2\par
State the degree of the polynomial.\par
0 0\par
\bigskip\textbf{Question 8.} Given that P(x) = Q(x)R (x) - Q (x)R(x), write down an expression for\par
P(x)\par
dx .\par
Z Q(x)\par
¡ ¢\par
\textbf{(i)} By choosing the function R(x) to be of the form a + bx + cx2, find\par
\[5x2 - 4x - 3\]
dx .\par
\[Z (1 + 2x + 3x2)2\]
Show that the choice of R(x) is not unique and, by comparing the two functions R(x)\par
corresponding to two different values of a, explain how the different choices are related.\par
\textbf{(ii)} Find the general solution of\par
dy\par
\[(1 + cos x + 2 sin x) + (sin x - 2 cos x)y = 5 - 3 cos x + 4 sin x .\]
dx\par
\[9475 Jun10 [Turn over\]
\end{flushleft}
\subsection*{Page 5}
\begin{flushleft}
\bigskip\textbf{Section B: Mechanics}\par
P\par
Q\par
a\par
θ\par
O\par
The diagram shows two particles, P and Q, connected by a light inextensible string which\par
passes over a smooth block fixed to a horizontal table. The cross-section of the block is a\par
quarter circle with centre O, which is at the edge of the table, and radius a. The angle\par
between OP and the table is θ. The masses of P and Q are m and M , respectively, where\par
\[m < M .\]
Initially, P is held at rest on the table and in contact with the block, Q is vertically above O,\par
and the string is taut. Then P is released. Given that, in the subsequent motion, P remains\par
in contact with the block as θ increases from 0 to \pi, find an expression, in terms of m, M ,\par
θ and g, for the normal reaction of the block on P and show that\par
\[m \pi - 1\]
\[> .\]
M 3\par
\bigskip\textbf{Question 10.} A small bead B, of mass m, slides without friction on a fixed horizontal ring of radius a. The\par
centre of the ring is at O. The bead is attached by a light elastic string to a fixed point P in\par
the plane of the ring such that OP = b, where b > a. The natural length of the elastic string\par
is c, where c < b - a, and its modulus of elasticity is λ. Show that the equation of motion of\par
the bead is\par
a sin φ\par
\[maφ¨ = -λ - 1 sin(θ + φ) ,\]
µ c sin θ ¶\par
\[where θ = ∠BP O and φ = ∠BOP .\]
Given that θ and φ are small, show that a(θ + φ) \approx bθ. Hence find the period of small\par
\[oscillations about the equilibrium position θ = φ = 0.\]
\[9475 Jun10\]
\end{flushleft}
\subsection*{Page 6}
\begin{flushleft}
\bigskip\textbf{Question 11.} A bullet of mass m is fired horizontally with speed u into a wooden block of mass M at rest\par
on a horizontal surface. The coefficient of friction between the block and the surface is µ.\par
While the bullet is moving through the block, it experiences a constant force of resistance to\par
\[its motion of magnitude R, where R > (M + m)µg. The bullet moves horizontally in the block\]
and does not emerge from the other side of the block.\par
\textbf{(i)} Show that the magnitude, a, of the deceleration of the bullet relative to the block while\par
the bullet is moving through the block is given by\par
R R - (M + m)µg\par
\[a = + .\]
m M\par
\textbf{(ii)} Show that the common speed, v, of the block and bullet when the bullet stops moving\par
through the block satisfies\par
Ru - (M + m)µgu\par
\[av = .\]
M\par
\textbf{(iii)} Obtain an expression, in terms of u, v and a, for the distance moved by the block while\par
the bullet is moving through the block.\par
\textbf{(iv)} Show that the total distance moved by the block is\par
muv\par
.\par
2(M + m)µg\par
\[Describe briefly what happens if R < (M + m)µg.\]
\[9475 Jun10 [Turn over\]
\end{flushleft}
\subsection*{Page 7}
\begin{flushleft}
\bigskip\textbf{Section C: Probability and Statistics}\par
\bigskip\textbf{Question 12.} The infinite series S is given by\par
\[S = 1 + (1 + d)r + (1 + 2d)r 2 + \cdots + (1 + nd)rn + \cdots ,\]
\[for |r| < 1 . By considering S - rS, or otherwise, prove that\]
1 rd\par
\[S = + .\]
1 - r (1 - r)2\par
Arthur and Boadicea shoot arrows at a target. The probability that an arrow shot by Arthur\par
hits the target is a; the probability that an arrow shot by Boadicea hits the target is b. Each\par
shot is independent of all others. Prove that the expected number of shots it takes Arthur to\par
hit the target is 1/a.\par
Arthur and Boadicea now have a contest. They take alternate shots, with Arthur going first.\par
The winner is the one who hits the target first. The probability that Arthur wins the contest\par
is α and the probability that Boadicea wins is β. Show that\par
a\par
\[α = ,\]
\[1 - a0b0\]
0 0\par
\[where a = 1 - a and b = 1 - b, and find β.\]
α β\par
Show that the expected number of shots in the contest is + .\par
a b\par
\bigskip\textbf{Question 13.} In this question, Corr(U, V ) denotes the product moment correlation coefficient between the\par
random variables U and V , defined by\par
Cov(U, V )\par
Corr(U, V ) \equiv .\par
Var (U )Var (V )\par
p\par
\[The independent random variables Z^1, Z^2 and Z^3 each have expectation 0 and variance 1.\]
\[What is the value of Corr(Z^1, Z^2)?\]
\[Let Y^1 = Z^1 and let\]
\[Y^2 = ρ\]
\[Z^1 + (1 - ρ 2\]
\[)2 Z^2 ,\]
where ρ\par
\[is a given constant with -1 < ρ\]
\[< 1. Find E(Y^2), Var (Y^2) and Corr(Y^1, Y^2).\]
\[Now let Y^3 = aZ1 + bZ2 + cZ3, where a, b and c are real constants and c > 0. Given that\]
\[E(Y^3) = 0, Var (Y^3) = 1, Corr(Y^1, Y^3) = ρ\]
\[and Corr(Y^2, Y^3) = ρ\]
, express a, b and c in\par
terms of ρ , ρ and ρ .\par
23 13 12\par
Given constants µ and σ , for i = 1, 2 and 3, give expressions in terms of the Y for random\par
i i i\par
\[variables X such that E(X ) = µ , Var (X ) = σ2 and Corr(X , X ) = ρ .\]
i i i i i i j ij\par
\[9475 Jun10\]
\end{flushleft}
\end{document}
